% ==============================================================================
% SECCION 1: METODOLOGIA Y ENTORNO DE SIMULACION FISICA
% ==============================================================================
\section{Metodología y Configuración del Entorno de Simulación}
\label{sec:metodologia}

Para el desarrollo del presente informe se implementó una infraestructura de simulación automatizada en SPICE y post-procesamiento en Python, orientada a garantizar reproducibilidad determinística y trazabilidad física absoluta en cada ensayo.

\subsection{Modelos de Dispositivos y Nodo Tecnológico}
Se emplean tarjetas de modelo predictivas BSIM4 (\textit{Predictive Technology Model}, PTM) desarrolladas por la Universidad Estatal de Arizona (ASU / Nanoscale Integration and Modeling Group). Con el propósito de prevenir colisiones de nombres al incluir múltiples nodos en simulaciones de referencia cruzada, los modelos se integraron en librerías modulares:
\begin{itemize}[leftmargin=1.4em,itemsep=0.1em]
  \item \texttt{ptm45hp.lib}: Proceso PTM de 45\,nm \textit{High Performance} ($V_{DD} = 1.0\,\text{V}$, $L_{\text{nom}} = 45\,\text{nm}$, $t_{oxe} = 1.25\,\text{nm}$, $V_{th0,n} \approx 0.47\,\text{V}$, $V_{th0,p} \approx -0.42\,\text{V}$). Nodo primario en la Actividad~1.
  \item \texttt{ptm45lp.lib}: Proceso PTM de 45\,nm \textit{Low Power} ($V_{DD} = 1.1\,\text{V}$, $L_{\text{nom}} = 45\,\text{nm}$, $t_{oxe} = 1.80\,\text{nm}$, $V_{th0,n} \approx 0.62\,\text{V}$, $V_{th0,p} \approx -0.58\,\text{V}$).
  \item \texttt{ptm130.lib}: Proceso clásico de 130\,nm Bulk ($V_{DD} = 1.3\,\text{V}$, $L_{\text{nom}} = 130\,\text{nm}$, $t_{oxe} = 3.10\,\text{nm}$).
\end{itemize}

\subsection{Parámetros Numéricos del Solver SPICE}
De acuerdo con las directrices de diseño submicrométrico para lógica dinámica, se estableció como directiva mandatoria en todos los archivos de netlist (\texttt{.cir}):
\begin{lstlisting}[style=spice]
.options gmin=1e-15 abstol=1e-14 reltol=1e-4 method=gear
\end{lstlisting}

La selección de estos parámetros responde a requerimientos físicos concretos:
\begin{enumerate}[leftmargin=1.4em,itemsep=0.12em]
  \item \textbf{Conductancia de regularización (\texttt{gmin = 1e-15}):} El valor predeterminado de LTspice (\texttt{gmin = 1e-12}, correspondiente a $1\,\text{pS}$) conecta una conductancia parásita hacia tierra en cada nodo, inyectando una corriente artificial calculable como $I_{gmin} = G_{min} \cdot V_{DD} = 10^{-12} \times 1.0\,\text{V} = 1.0\,\text{pA}$. En transistores de 45\,nm en estado de reposo, esta corriente numérica es comparable a las fugas físicas subumbral, distorsionando severamente las evaluaciones de retención de carga en nodos flotantes. Reducir \texttt{gmin} a $1\,\text{fS}$ minimiza esta inyección espuria a $1.0\,\text{aA}$ ($10^{-18}\,\text{A}$).
  \item \textbf{Tolerancias de convergencia (\texttt{abstol = 1e-14}, \texttt{reltol = 1e-4}):} Las corrientes residuales subumbral y de túnel en dispositivos en corte se ubican en el rango de femtoamperios. La tolerancia estándar ($10^{-12}\,\text{A}$) induce ruido numérico y falsas oscilaciones durante la evaluación de nodos dinámicos.
  \item \textbf{Algoritmo de integración (\texttt{method = gear}):} El método trapezoidal estándar (\texttt{trap}) produce resonancias numéricas artificiales (\textit{trapezoidal ringing}) cuando ocurren transiciones de tensión con tiempos de subida extremadamente reducidos ($t_r = 20\,\text{ps}$) sobre nodos con capacitancias parásitas de femtofaradios. La integración de Gear proporciona un amortiguamiento numérico incondicional sin alterar la constante de tiempo física $RC$.
\end{enumerate}

\subsection{Flujo de Extracción y Trazabilidad}
Todas las simulaciones fueron procesadas en modo \textit{batch} silencioso (\texttt{-b}) mediante LTspice v26.0.2 bajo Windows. Las figuras fueron generadas directamente a partir de los datos crudos extraídos de los ficheros \texttt{.raw} a 300\,DPI en formato matricial con ejes normalizados, mientras que las métricas escalares fueron recolectadas de los ficheros \texttt{.log} mediante directivas \texttt{.meas} acotadas en tiempo y un parser automatizado de regresión.
